\documentclass{article}

% NeurIPS 2025 style (vendored in neurips/neurips_2025.sty), preprint mode:
% these are self-study reports, not conference submissions.
\PassOptionsToPackage{numbers,sort&compress}{natbib}
\usepackage[preprint]{neurips/neurips_2025}

\usepackage[T1]{fontenc}
\usepackage{amsmath,amssymb}
\usepackage{booktabs}
\usepackage{graphicx}
\usepackage{array}
\usepackage[colorlinks=true,linkcolor=blue,urlcolor=blue,citecolor=blue]{hyperref}

\newcommand{\codex}[0]{Codex}
\newcommand{\occ}[0]{OpenCode}
\newcommand{\ccc}[0]{Claude Code}

\title{How Claude Code, Codex, and OpenCode Build Their Coding-Agent Harnesses:\\
A Static Source-Trace Comparison}
\author{self-study-with-ai}
\date{\today}

\begin{document}
\maketitle

\begin{abstract}
We compare how three production coding agents, Claude Code, Codex, and
OpenCode, implement the harness that turns a language model into a coding
agent: the turn loop, the tool surface, context compaction, permission and
sandbox policy, extensibility, configuration, session state, and interfaces.
The Claude Code core is closed source, so we pin all three repositories,
restrict the study to static source traces and official documentation
snapshots, and pre-register four trace experiments. Every extracted constant
and structure carries a file:line or snapshot-line anchor, and 300 anchors
re-verify against the pinned commits in a validation pass. The three systems
realize the same abstract turn loop with structurally different control
flow: an abortable task state machine in \codex{}, a persistence-driven
while loop in \occ{}, and docs-attested hook cadences in \ccc{}. Patch
application converges, with \occ{} porting the \codex{} \path{apply_patch}
format and selecting it by model ID. Safety postures diverge structurally:
\codex{} pairs an OS sandbox with a Starlark rule engine, \occ{} runs an
in-process ruleset with no OS sandbox in the pinned tree, and \ccc{}
documents an OS sandbox scoped to Bash commands only. Compaction triggers
are literal constants of different kinds, a fraction of the context window
in \codex{} against a reserved buffer in \occ{}, while \ccc{} discloses no
numeric threshold. All findings are bounded to the pinned commits
(\path{af70018}, \path{d545d8f}, \path{c3d2e35}) and the documentation
snapshots of 2026-08-20.
\end{abstract}

\section{Introduction}

A coding agent is a loop: the model proposes, tools act on a repository,
observations return, and the harness decides whether to continue. This
decomposition has research pedigree. ReAct interleaves reasoning traces
with environment actions and showed the combination beats either alone on
question answering and interactive tasks~\citep{yao2023react}. Toolformer
trained a model to emit, execute, and resume around API calls,
establishing the interrupt-execute-resume protocol that every agent harness
now re-implements in engineering form~\citep{schick2023toolformer}. SWE-agent
demonstrated that the interface itself is a first-order factor: an
agent-tailored computer interface with a guarded edit command and lint
gating raised SWE-bench Lite resolution from 11.0\% (shell-only) to 18.0\%
with GPT-4 Turbo, without changing the model~\citep{yang2024sweagent}.

What the literature leaves open is how this interface layer is built in
the shipping products that practitioners use. We study three of
them, \ccc{}, \codex{}, and \occ{}, at pinned source versions, and ask four
questions: what components make up a coding-agent harness; where the
three systems genuinely differ; how each trades capability against safety;
and what the closed \ccc{} core reveals through documentation, its plugin
surface, and third-party teardowns.

The \ccc{} core is closed source, and its pinned public checkout contains no
harness code~\citep{claudeCodeSurface}. We therefore treat documentation
snapshots as the evidence tier for that system, hold third-party teardowns
to hedged context only, and mark closed-core internals as boundaries rather
than findings. Contributions at the supported scope:

\begin{itemize}
  \item an eight-dimension decomposition of a coding-agent harness, grounded
  in three pinned implementations and the framing literature
  (Section~\ref{sec:components});
  \item four pre-registered static-trace extractions, turn loops, tool
  surfaces, safety policies, and compaction parameters, whose anchors all
  re-verify at the pinned commits (Section~\ref{sec:methods},
  Sections~\ref{sec:turnloop} to~\ref{sec:compaction});
  \item a dimension-by-dimension comparison that separates genuine
  structural differences from naming differences
  (Sections~\ref{sec:extensibility} to~\ref{sec:convergence}).
\end{itemize}

\section{Background}
\label{sec:background}

\paragraph{The abstract loop.} ReAct frames agency as an interleave of
thoughts, actions, and observations, and measured both its gains and its
failure modes: grounded search cut hallucinated failures from 56\% (chain
of thought) to 0\% of failures on HotpotQA, at the cost of a repetitive
failure pattern in which the model re-generates its previous thoughts and
actions~\citep{yao2023react}. Harness designers inherit both sides: a
continuation condition, and a defense against loops.

\paragraph{Tool calling as protocol.} Toolformer learns when to call a tool,
with what arguments, and how to fold results back into prediction, and
restricts decoding to at most one API call per input explicitly to prevent
runaway tool loops~\citep{schick2023toolformer}. The same concern appears in
all three harnesses as explicit machinery (Section~\ref{sec:turnloop}).

\paragraph{The agent-computer interface.} SWE-agent ablated interface
choices on SWE-bench Lite and found each one measurable: removing the edit
command dropped resolution by 7.7 points, disabling linting dropped it by
3.0 points, and an iterative search interface performed worse than no added
search tools at all~\citep{yang2024sweagent}. Editing is where agents spend
their failures: 51.7\% of instances had at least one failed edit, and the
probability of eventual edit success dropped from 90.5\% to 57.2\% after a
single failure~\citep{yang2024sweagent}. Guardrails trade flexibility for
recovery, a tradeoff our three systems resolve differently
(Sections~\ref{sec:tools},~\ref{sec:safety}).

\section{Methods and Sources}
\label{sec:methods}

\paragraph{Design.} Descriptive static source-trace study. The units of
analysis are literal code constructs at pinned commits and quoted
documentation text at pinned snapshots. There is no behavioral measurement:
no agent CLI was invoked, no live API was called, and no code from the
pinned repositories was executed. The pinned checkouts were treated as
read-only untrusted data.

\paragraph{Pins and sources.} Three local checkouts were pinned on
2026-08-19 at commits \path{af700180808cce2ce28a31aad0fbad4dc58b857a} (\codex{},
Apache-2.0), \path{d545d8fba57283528db69281f59c803c646eb7e9} (\occ{}, MIT,
\texttt{dev} branch), and \path{c3d2e35e554060b5a20ee6b28140fbdbd4eb0048}
(\ccc{})~\citep{codexRepo,opencodeRepo,claudeCodeSurface}. Official
documentation snapshots were fetched 2026-08-20; documentation sites are
not version-pinned and can drift, so documentation claims cite the snapshot
rather than a commit~\citep{claudeCodeHooksDocs,claudeCodeSubagentsDocs,%
claudeCodePluginsDocs,claudeCodeMcpDocs,claudeCodeMemoryDocs,%
claudeCodePermissionsDocs,claudeCodeSandboxingDocs,codexSandboxingDocs,%
opencodeDocs,codexDocsRepo}. Evidence tiers: pinned code (primary for \codex{}
and \occ{}), pinned plugin/example surface plus documentation snapshots
(only tier available for the \ccc{} core), framing preprints
(background), and third-party teardowns, \texttt{tier: blog}, admissible
only for hedged contextual statements~\citep{minusXTeardown,agiflowInternals,%
tenguDecoded}.

\paragraph{Trace experiments.} Four experiments were frozen in a plan
before any full extraction (EXP-PLAN-2026-08-19-v1): E1 turns loops, E2
tool surfaces, E3 safety policies, E4 compaction parameters. Each produced
markdown and JSON artifacts in which every stated value carries a
\path{file:line} or snapshot-line anchor, plus a validation script that
re-checks every anchor against the pinned sources. Results: 59/59 anchor
checks for E1, 131/131 for E2, 75/75 for E3, 35/35 for E4
(300 total); all three HEADs matched the pins at every validation. The
validation runs are smoke-test plumbing, not scientific evidence; the
extraction artifacts carry the findings. Twelve claims advanced to
ANALYZED with supported evidential status (\texttt{CLM-E1-1a} through
\texttt{CLM-E4-3a} in the study's claim ledger, each superseding its
EXECUTED counterpart; after review, \texttt{CLM-E2-1a} was itself
superseded by \texttt{CLM-E2-1b} correcting a registration count from 33
to 32).

\paragraph{Closed-core protocol.} The \ccc{} column is filled exclusively
from documentation quotes and the pinned plugin/example surface. Where no
pinned source states a value, the text says so explicitly and the gap is
listed in Section~\ref{sec:limitations}. Blog-tier teardowns never fill a
comparison cell on their own~\citep{claudeCodeSurface,agiflowInternals}.

\section{Findings}

\subsection{What makes a coding-agent harness}
\label{sec:components}

Cross-system, the harness decomposes into eight machinery dimensions, each
independently implemented by all three systems: the turn loop; the tool
surface and its file-edit protocol; context assembly and compaction;
persistent memory files; permission and sandbox policy; extensibility
(hooks, plugins, skills, subagents, MCP servers); configuration and model
provider plumbing; and session state with durable storage and user-facing
transports. Table~\ref{tab:summary} summarizes the matrix; the details
follow per dimension.

\begin{table}[t]
\centering
\small
\caption{Summary comparison. Anchors: E1 to E4 are the trace experiment
artifacts; notes are the per-component study notes
(Section~\ref{sec:methods}). ``Undisclosed'' means no pinned source states
the value.}
\label{tab:summary}
\begin{tabular}{p{2.5cm}p{4.1cm}p{4.1cm}p{4.1cm}}
\toprule
Dimension & \codex{} (\path{af70018}) & \occ{} (\path{d545d8f}) & \ccc{} (docs + surface)\\
\midrule
Turn loop & abortable task state machine; 27-op dispatch; continuation bit \path{needs_follow_up} (E1, \path{codexTurnLoop}) & persistence-driven \path{while (true)} loop; step verdict \path{compact|stop|continue} (E1, \path{opencodeSessionLoop}) & closed core; hook cadences per session, per turn, per tool call (E1, hooks snapshot)\\
Tools & 32 conditional registrations; \path{apply_patch} grammar; no read/glob/grep tools (E2, \path{codexToolsPatch}) & fixed 17-ID array; \path{apply_patch} port gated on \path{gpt-*} model IDs (E2, \path{opencodeTools}) & 12 built-ins named in docs; \path{ToolSearch} default-on (E2, hooks and MCP snapshots)\\
Compaction & fraction trigger, 9/10 of 95\% window; explicit retention budgets (E4, \path{codexContextCompaction}) & reserved-buffer trigger, input limit minus a reservation defaulting to min(20{,}000, \path{maxOutputTokens}) tokens (E4, \path{opencodeContextCompaction}) & seams only; thresholds undisclosed (E4, hooks and memory snapshots)\\
Safety & OS sandbox plus Starlark execpolicy, strictest-wins (E3, \path{codexSandboxPermissions}) & in-process ruleset only; no OS sandbox in tree (E3, \path{opencodePermissions}) & permission rules plus Bash-only OS sandbox (E3, permissions and sandboxing snapshots)\\
Extensibility & hooks with content-hash trust; plugins as \path{<plugin>@<marketplace>} (\path{codexExtensibility}) & in-process TypeScript plugins; 7 built-in agents (\path{opencodeExtensibility}) & 31 hook events; 10-component plugins; subagent quotas (hooks, plugins, subagents docs)\\
Config & 10 numeric precedence layers; trust-gated project config (\path{codexConfigProviders}) & 9 layers ending in MDM; runtime provider install (\path{opencodeConfigProviders}) & closed loader; managed settings documented as enforcement (memory, MCP docs)\\
State & append-only JSONL plus repairable SQLite mirror (\path{codexStateRollout}) & event-sourced SQLite plus shadow-git snapshots (\path{opencodeStateSnapshots}) & undisclosed; subagent transcripts documented (subagents docs)\\
Interfaces & every frontend is one JSON-RPC protocol client (\path{codexInterfaces}) & TUI consumes local HTTP server; ACP; 38 LSP servers (\path{opencodeInterfaces}) & CLI, IDE extensions, desktop app attested in docs (permissions docs)\\
\bottomrule
\end{tabular}
\end{table}

\subsection{Turn loops}
\label{sec:turnloop}

All three harnesses implement the same abstract loop and realize it
structurally differently (Table~\ref{tab:turnloop}).

\begin{table}[t]
\centering
\small
\caption{Turn-loop structures at the pinned commits. \ccc{} values are
docs-attested; its implementation values are undisclosed (closed core).}
\label{tab:turnloop}
\begin{tabular}{p{2.6cm}p{4.0cm}p{4.0cm}p{4.0cm}}
\toprule
Property & \codex{} & \occ{} & \ccc{} (docs)\\
\midrule
Realization & background Tokio task; \path{Submission} envelope; \path{Op} enum, 27 variants; loop exits on \path{Op::Shutdown} (\path{protocol.rs:543-698}) & \path{while (true)} over persisted message parts, \path{prompt.ts:1088} & closed; per-turn \path{UserPromptSubmit ... Stop/StopFailure} nesting per-tool-call \path{PreToolUse/PostToolUse} (hooks snapshot :20-24)\\
Continuation & \path{needs_follow_up = model_needs_follow_up || has_pending_input}, \path{session/turn.rs:405} & exit needs non-\path{tool-calls} finish, no pending tool parts, assistant parented to last user message (\path{prompt.ts:1106-1130}) & undisclosed\\
Step verdict & n/a (task-level: 3 \path{TaskKind} variants Regular, Review, Compact) & \path{compact|stop|continue}, \path{processor.ts:30} & undisclosed\\
Abort & 4 \path{TurnAbortReason} variants; spawn preempts with \path{Replaced} (\path{tasks/mod.rs:279-288}) & abort routes; denied permissions break the loop unless an experimental flag is set & undisclosed\\
Anti-loop & 8-variant \path{NotSubmittedReason} steering rejections (\path{turn_input.rs:182-208}) & \path{DOOM_LOOP_THRESHOLD = 3} byte-identical calls raise a permission ask (\path{processor.ts:29}) & stop hooks capped after 8 consecutive blocks (hooks snapshot :2397)\\
Retry and caps & n/a at loop level & \path{RETRY_MAX_RETRIES = 5}, initial delay 2000 ms, backoff factor 2 (\path{retry.ts:26-31}); per-agent \path{steps} cap is a soft guardrail & undisclosed; docs cap agent hooks at 50 turns (hooks snapshot)\\
\bottomrule
\end{tabular}
\end{table}

\codex{} models a turn as an abortable task inside a state machine: inputs
arrive as \path{Submission} envelopes whose \path{Op} enum has exactly 27
variants at the pinned commit; any new task preempts a running one with
\path{TurnAbortReason::Replaced}; and the sampling loop continues exactly
while \path{needs_follow_up = model_needs_follow_up || has_pending_input},
where the model-side bit is set by queued tool calls, by tool errors marked
\path{RespondToModel}, and by completed responses that did not end the
turn~\citep{codexRepo}. Mid-turn, a rollover gate runs automatic compaction
before continuing (Section~\ref{sec:compaction}). Steering an active turn is
itself a typed decision with an 8-variant rejection taxonomy; accepted
replies mean only that the core accepted the input for processing, not that
it will act on it~\citep{codexRepo}.

\occ{} is persistence-driven: the loop condition is recomputed each
iteration from stored message parts, not from in-memory step state. Exiting
requires all of a non-\path{tool-calls} finish reason, no pending tool
parts, and the last assistant message parented to the last user message; the
exit-condition comment explains the redundancy as defense against providers
that return \path{stop} while the message still contains tool
calls~\citep{opencodeRepo}. Each step inside the loop returns one of three
verdicts, \path{compact}, \path{stop}, or \path{continue}; a repeated
identical tool call (three times, byte-identical serialized input) raises a
\path{doom_loop} permission ask whose default action is itself
\path{ask}~\citep{opencodeRepo}.

\ccc{} exposes no loop code. The hook reference documents three event
cadences that attest the loop's seam structure: once per session, once per
turn, and on every tool call inside the agentic loop, and states that a
blocking stop hook is overridden after 8 consecutive blocks~\citep{%
claudeCodeHooksDocs}. The continuation condition, any step cap, and the
retry policy are not documented in the pinned snapshots
(Section~\ref{sec:limitations}).

\subsection{Tool surfaces and patch application}
\label{sec:tools}

The tool shapes overlap (one shell path, at least one edit path, search,
delegation), but the registration policy and the edit protocol differ
(Table~\ref{tab:tools}).

\begin{table}[t]
\centering
\small
\caption{Tool surfaces. \codex{} and \occ{} cells are code-anchored; \ccc{}
cells are docs-attested or attested by the pinned plugin surface (closed
core). Undisclosed cells have no pinned source.}
\label{tab:tools}
\begin{tabular}{p{2.2cm}p{4.2cm}p{4.2cm}p{4.2cm}}
\toprule
Shape & \codex{} & \occ{} & \ccc{} (docs/surface)\\
\midrule
Shell & \path{shell_command} (legacy) or \path{exec_command} + \path{write_stdin} (unified pty), chosen by shell tool type & \path{bash} & \path{Bash}, plus \path{PowerShell} (hooks snapshot :1513)\\
File edit & \path{apply_patch}, freeform string parsed by a 19-line Lark grammar & \path{edit} (9-stage replacer cascade, similarity 0.65) and \path{write}; \path{apply_patch} port & \path{Edit}, \path{Write}, \path{MultiEdit}, \path{NotebookEdit} (surface attestations)\\
Patch-vs-edit choice & gated on \path{model_info.apply_patch_tool_type} & \path{apply_patch} replaces \path{edit}/\path{write} for \path{gpt-*} except \path{gpt-4}/\path{oss} (\path{registry.ts:296-300}) & undisclosed\\
Read/search & no dedicated tools; shell carries read and search; \path{view_image} for images & \path{read}, \path{glob}, \path{grep} & \path{Read}, \path{Glob}, \path{Grep} (hooks snapshot :1513)\\
Delegation & 6 v2 agent tools (\path{spawn_agent} ... \path{list_agents}) or 5 v1 tools, feature-gated & \path{task} & \path{Agent} (docs), \path{Task} (surface); \path{SubagentStart/Stop} hook events\\
Planning & \path{update_plan} behind config & \path{plan_exit} behind experimental flags & \path{ExitPlanMode} (hooks snapshot :1513)\\
Discovery & \path{tool_search} registered only when a deferred searchable tool exists & none builtin & \path{ToolSearch} enabled by default; doubles as the MCP wait point (MCP snapshot :279,:310)\\
Context tooling & \path{new_context} and \path{get_context_remaining} behind a feature gate & session-managed compaction, no model-visible tool & no model-visible tool attested; \path{PreCompact/PostCompact} hook seams\\
Shell defaults & 10 s timeout, 1 MiB output cap, 10{,}000 event deltas per call; unified exec: 10 s yield default clamped at 30 s, 300 s background, 64 processes & truncation via \path{tool/truncate.ts}; numeric caps outside E2 scope & Bash sandbox scope is Bash-only (surface); numeric defaults undisclosed\\
\bottomrule
\end{tabular}
\end{table}

\codex{} exposes the widest conditional surface. Tool registration happens
per turn: 32 handler-to-tool-name mappings, each with an explicit condition
(feature gates, model capabilities, or the presence of other
tools)~\citep{codexRepo}. \path{apply_patch} is a grammar-constrained
freeform tool, parsed by a 19-line Lark grammar, and is registered only when
the model's \path{apply_patch_tool_type} says so. \codex{} has no dedicated
read, glob, or grep tool in the extracted registry; file inspection is
delegated to the shell tools~\citep{codexRepo}. Shell defaults are
explicit: a 10 s timeout, a 1 MiB output cap, and 10{,}000 output
event deltas per call.

\occ{} ships a fixed ordered array of 17 tool IDs and makes one structural
swap: for model IDs containing \path{gpt-}, except \path{gpt-4} and
\path{oss}, \path{apply_patch} replaces \path{edit} and
\path{write}~\citep{opencodeRepo}. The swap exists because \occ{} ports the
\codex{} patch format directly; the port's comment reads ``Core types
matching the Rust implementation'' (\path{patch/index.ts:12}). The
\path{edit} tool itself falls back through an ordered 9-stage replacer
cascade with Levenshtein similarity thresholds of 0.65, an engineering
response to the same fragility SWE-agent measured, where edit success drops
sharply after one failed edit~\citep{yang2024sweagent,opencodeRepo}.

\ccc{}'s inventory is attestable only through documentation and the pinned
surface. The hooks snapshot names 12 built-in tools plus MCP tools; the
pinned plugin and example surface attests 8 more names (LS, NotebookRead,
TodoWrite, KillShell, BashOutput, Task, MultiEdit, NotebookEdit), and the
complete closed-core registry remains open~\citep{claudeCodeHooksDocs,%
claudeCodeSurface}. Tool discovery is implemented differently where it
exists: \codex{} registers \path{tool_search} only when a deferred
searchable tool exists, while \ccc{} documents \path{ToolSearch} as enabled
by default and as the point where the harness waits for MCP
connections~\citep{codexRepo,claudeCodeMcpDocs}. Third-party teardowns
assert a larger inventory (49 tools in one binary analysis) and a tool-tier
design (low, medium, high level), but these are unverified blog-tier
observations~\citep{tenguDecoded,minusXTeardown}.

\subsection{Context compaction}
\label{sec:compaction}

Both open systems trigger compaction at literal constants, and the
constants are of different kinds (Table~\ref{tab:compaction}).

\begin{table}[t]
\centering
\small
\caption{Compaction parameters at the pinned commits. \ccc{} exposes the
compaction lifecycle at hook seams but discloses no numeric thresholds in
the pinned snapshot set.}
\label{tab:compaction}
\begin{tabular}{p{2.6cm}p{4.0cm}p{4.0cm}p{4.0cm}}
\toprule
Parameter & \codex{} & \occ{} & \ccc{} (docs)\\
\midrule
Trigger style & fractional & reserved buffer & undisclosed\\
Threshold & 9/10 of the resolved context window, capped by config (\path{openai_models.rs:488-499}) & count reaches input limit minus a reservation defaulting to min(20{,}000, \path{maxOutputTokens}) tokens (\path{overflow.ts:8-34}) & undisclosed\\
Headroom & 95\% effective window default (\path{openai_models.rs:376-378}) & \texttt{COMPACTION\_BUFFER} = 20{,}000 tokens & undisclosed\\
Summarization budget & user messages capped at 20{,}000 tokens locally (\path{compact.rs:57}) & pruning targets tool outputs, not user messages & undisclosed; root CLAUDE.md re-injected after \path{/compact}, nested files are not (memory snapshot :450-454)\\
Retention & remote path retains 64{,}000 tokens, agent messages capped at 10{,}000, 2 retries (\path{compact_remote_v2.rs:65-69}) & preserve clamp(25\% of usable, 2{,}000 to 15{,}000) tokens of recent messages (\path{compaction.ts:115-118}) & undisclosed\\
Pruning & summary replaces context; rollback op drops last N turns & keep last 40{,}000 tokens of tool output (\path{PRUNE_PROTECT}); minimum 20{,}000 before pruning; truncated outputs keep 2{,}000 chars; skill outputs never pruned & \path{PostCompact} input carries \path{compact_summary} (hooks docs)\\
Mid-turn & rollover runs auto-compaction then continues (ContextLimit, MidTurn) & verdict \path{compact} appends a compaction part and continues (\path{processor.ts:679}) & blocking \path{PreCompact} during context-limit recovery surfaces the error (hooks docs)\\
Token accounting & model-reported counts & chars/4 estimate for unknown counts (\path{token.ts:3-5}) & undisclosed\\
\bottomrule
\end{tabular}
\end{table}

\codex{} computes \path{auto_compact_token_limit} as nine tenths of the
resolved context window (capped by any explicit configuration limit), on
top of an effective window held at 95\% by default; local summarization
caps the user messages fed to the summarizer at 20{,}000 tokens, and the
remote summarization path retains a 64{,}000-token budget with
10{,}000-token agent-message caps and 2 retries~\citep{codexRepo}. Mid-turn,
the rollover gate runs auto-compaction and continues the loop, so
compaction is a state transition inside the turn rather than between turns.

\occ{} triggers on an absolute boundary: overflow fires when the running
token count reaches the input limit minus a reserved buffer defaulting to
min(20{,}000, \path{maxOutputTokens}) tokens~\citep{opencodeRepo}. Its
budgets are pruning-shaped rather than summary-shaped: the pruner protects
the last 40{,}000 tokens of tool output, requires at least 20{,}000 tokens
before it runs, truncates outputs to 2{,}000 characters, never prunes skill
tool outputs, and preserves recent messages clamped to 25\% of usable
context within 2{,}000 to 15{,}000 tokens. Unknown token counts are
estimated as characters divided by four.

\ccc{} documents the lifecycle but not the numbers. The hook reference
carries \path{PreCompact}/\path{PostCompact} events with \path{manual}/%
\path{auto} trigger values, and states that a hook blocking compaction
during context-limit recovery surfaces the underlying error; the memory
documentation states that project-root CLAUDE.md survives \path{/compact}
by being re-read and re-injected, while nested CLAUDE.md files do
not~\citep{claudeCodeHooksDocs,claudeCodeMemoryDocs}. No numeric trigger,
summarization budget, or retention policy appears anywhere in the pinned
snapshot set (Section~\ref{sec:limitations}).

\subsection{Permissions and sandboxing}
\label{sec:safety}

The three systems occupy three distinct points on the same axis
(Table~\ref{tab:safety}).

\begin{table}[t]
\centering
\small
\caption{Safety policies. \ccc{} cells cite the pinned documentation
snapshots and the pinned settings surface; enforcement internals are closed.
Undisclosed cells have no pinned source.}
\label{tab:safety}
\begin{tabular}{p{2.6cm}p{4.0cm}p{4.0cm}p{4.0cm}}
\toprule
Dimension & \codex{} & \occ{} & \ccc{} (docs)\\
\midrule
Policy engine & Starlark rule engine over \path{.codexpolicy} files & in-process permission ruleset & permission rules plus OS sandbox as two documented layers (sandboxing snapshot :521-525)\\
Decision model & \path{Decision = Allow|Prompt|Forbidden} (\path{decision.rs:9-16}) & \path{allow|ask|deny} per ruleset entry & six modes default, acceptEdits, plan, auto, dontAsk, bypassPermissions (permissions snapshot :56-63); rule order deny, ask, allow (:42)\\
Aggregation & strictest-wins via \path{max()} over matched rules (\path{policy.rs:402-403}) & merge precedence not extracted & hook decisions \path{deny > defer > ask > allow} (hooks docs)\\
Shipped content & no \path{.rules} file; only a 78-line example marked not recommended for actual use & default ruleset ships: \path{"*": allow}, doom loop ask, \path{question/plan_enter/plan_exit} deny, \path{*.env} ask (\path{agent.ts:119-136}) & rule syntax documented; no shipped default policy text in snapshots\\
OS sandbox & Seatbelt (macOS, 122-line compiled policy), bubblewrap argv plus seccomp/landlock network filtering (Linux) & none found in the pinned tree & built-in Seatbelt (macOS); bubblewrap plus socat (Linux, WSL2); native Windows not supported (sandboxing snapshot)\\
Sandbox scope & 4 \path{SandboxPolicy} variants; \path{WritableRoot} protects \path{.git}, \path{.agents}, \path{.codex} (\path{protocol/src/permissions.rs:24-33}) & n/a & Bash commands and their child processes only; Read/Edit/Write use the permission system directly (sandboxing snapshot :523,:658)\\
Network control & seccomp denies ptrace/\path{io_uring}; restricted mode denies sockets except \path{AF_UNIX} & n/a & no domains pre-allowed; first connection prompts; allow for the session (sandboxing docs)\\
Evasion resistance & 90-line banned list of amendable shell/runtime prefixes (\path{exec_policy.rs:58-147}) & bash arity normalization, longest-prefix-wins over about 130 entries, before rule matching & fixed wrapper stripping; argument-scoped allow patterns documented as fragile (permissions snapshot)\\
Kill switches & enterprise \path{requirements.toml}; trust gating of project config & environment flags & \path{disableBypassPermissionsMode}/\path{disableAutoMode} (permissions snapshot :69)\\
\bottomrule
\end{tabular}
\end{table}

\codex{} layers mechanisms: an OS sandbox per platform under a rule
engine~\citep{codexRepo}. The execpolicy engine evaluates shell commands
against Starlark rules with a three-valued decision and strictest-wins
aggregation; no policy file ships with the repository, and the only example
is explicitly labeled as not recommended for actual use, so the default
behavior is approval-driven rather than policy-driven. The sandbox side is
concrete: a 122-line Seatbelt base policy compiled into the binary on
macOS, a bubblewrap argument builder plus a seccomp filter that
unconditionally denies ptrace and \path{io_uring} socket families on
Linux, and writable-root structures that deny writes to \path{.git},
\path{.agents}, and \path{.codex} inside permitted roots. Hardening against
command-prefix smuggling is a static 90-line list of amendable shell and
runtime prefixes~\citep{codexRepo,codexSandboxingDocs}.

\occ{} takes the opposite position: all safety is in-process. The pinned
tree contains no OS sandbox machinery within the searched scope, and the
default ruleset is permissive, \path{"*": allow}, with narrow friction at
\path{.env} reads (ask), the doom-loop tripwire (ask), and external
directories (ask), plus denials for the interaction and plan tools. Its
answer to command smuggling is normalization: shell commands are reduced to
a canonical arity via a longest-prefix-wins dictionary of about 130 entries
before rules match~\citep{opencodeRepo,opencodeDocs}.

\ccc{} documents a two-layer model. Permission rules are evaluated before
any tool runs and apply to every tool, in the fixed order deny, then ask,
then allow, with hook-level decisions ordered \path{deny > defer > ask >
allow}; the OS sandbox ``applies only to Bash commands and their child
processes'', while Read, Edit, and Write ``use the permission system
directly''~\citep{claudeCodePermissionsDocs,claudeCodeSandboxingDocs,%
claudeCodeHooksDocs}. Enforcement is platform-native (Seatbelt on macOS,
bubblewrap plus socat on Linux and WSL2), no domain is pre-allowed at the
network layer, and protected paths such as \path{.claude} settings,
\path{.git/hooks}, and shell startup files deny writes even inside writable
roots, with no per-path exemption. Enterprise kill switches disable
\path{bypassPermissions} and \path{auto}
modes~\citep{claudeCodePermissionsDocs,claudeCodeSandboxingDocs}. The
pinned surface corroborates the Bash-only scoping and adds the reference
confinement for untrusted work: a devcontainer whose firewall defaults to
dropping egress outside an allowlist~\citep{claudeCodeSurface}. The
implementation of these layers is closed (Section~\ref{sec:limitations}).

\subsection{Extensibility}
\label{sec:extensibility}

All three expose hooks, plugins, skills, subagents or agents, and MCP
integration, but the extension substrate differs more than the vocabulary
(Table~\ref{tab:summary}, row Extensibility).

\codex{} implements eleven hook events with regex matchers, and treats hook
code as untrusted input: handlers are content-hashed and run only when
marked managed or trusted~\citep{codexRepo}. \path{PreToolUse} hooks can
block, inject context, or rewrite tool input, with exit code 2 as the block
signal. Plugins are \path{<plugin>@<marketplace>} packages bundling skills,
hooks, MCP configuration, and app manifests, discovered under
\path{.codex-plugin}, \path{.claude-plugin}, and \path{.cursor-plugin}
directories. The \ccc{} compatibility is explicit in code: plugin hook
environments export \path{CLAUDE_PLUGIN_ROOT} and \path{CLAUDE_PLUGIN_DATA}
``for OOTB compat with existing plugins''. Skills are \path{SKILL.md}
documents with a 2\% of context window budget (capped at 10{,}000 tokens),
and MCP servers are config-defined with per-step policy overlays and a
protocol version pinned to 2025-06-18 by default.

\occ{} makes plugins a first-class in-process runtime: a plugin is a
TypeScript function returning hooks, dispatched sequentially in registration
order, and ten of the eleven internal plugins are provider authentication
bundles that exercise the same API~\citep{opencodeRepo}. It defines seven
built-in agents, including dedicated \path{compaction}, \path{title}, and
\path{summary} agents, with \path{explore} read-only by permission ruleset
except for an explicit \path{bash: "allow"} (\path{agent.ts:205}). Its MCP
client advertises only the \path{roots} capability, with sampling,
elicitation, and tasks disabled. Skills are discovered from \occ{}
directories,
from the vendor-neutral \path{.agents} layout, and from the \ccc{}
\path{.claude} layouts unless disabled, and every skill doubles as a slash
command~\citep{opencodeRepo,opencodeDocs}. The trust surface is
correspondingly wide: plugins receive the full SDK client in-process, and
the SDK declares a \path{permission.ask} hook for which no trigger call
site exists at this commit.

\ccc{} documents the broadest lifecycle surface: 31 hook events with five
handler types (command, HTTP, MCP tool, prompt, agent), hooks running in
parallel, and agent hooks capped at 50 turns~\citep{claudeCodeHooksDocs}.
Plugins carry up to ten component kinds, including monitors whose stdout
lines are delivered as session notifications and \path{bin/} executables
added to the Bash tool's path; two public marketplaces distribute them, one
curated and one community-reviewed~\citep{claudeCodePluginsDocs}. Subagents
are quota-controlled by documentation: a depth of three layers below the
main conversation by default, a 20-concurrent cap, and forks that inherit
the full conversation while remaining unable to spawn further
forks~\citep{claudeCodeSubagentsDocs}. The pinned surface shows how far the
hook contract delegates: an unbounded agent loop is implemented entirely in
a shell script that blocks the Stop event repeatedly
(the ralph-wiggum plugin), with no core changes~\citep{claudeCodeSurface}.

\subsection{Configuration and providers}
\label{sec:config}

\codex{} ranks configuration provenance numerically across ten layers,
from packaged defaults (-10) through MDM-delivered legacy managed config
(50), merged by recursive table merge in which overlay layers replace
non-table values wholesale and arrays never union across
layers~\citep{codexRepo}. Two mechanisms couple configuration to safety:
project-local configuration is disabled until the project is trusted, and a
twelve-key denylist hard-blocks repository content from setting provider
endpoints, model providers, notify, profile, realtime, and otel
keys (\path{loader/mod.rs:71-84}), with the source comment that
repo-local config should not choose where credentials are sent. The
provider layer is deliberately narrow: exactly five built-in provider IDs
(\path{openai}, \path{amazon-bedrock}, \path{amazon-bedrock-runtime},
\path{ollama}, \path{lmstudio}), reserved against override except Bedrock,
and a wire protocol pinned to the Responses API, with \path{chat} rejected
as a hard error pointing to a migration discussion.

\occ{} merges a layered precedence chain that ends in admin-managed config
directories and a macOS MDM plist whose code comment states it overrides
everything, with one asymmetry: inside
\path{.opencode} directories the entry farther from the working directory
wins, the opposite of plain-file precedence~\citep{opencodeRepo}. Its
provider layer is the widest of the three: 24 bundled AI-SDK provider
packages, runtime installation of arbitrary provider SDKs named in config,
and a hosted model catalog fetched from models.opencode.ai and cached for
five minutes. Configuration handling is capability-first: unknown keys are
silently discarded, environment substitution falls back to empty without
error, and an inline environment variable replaces the credentials file
without schema validation.

\ccc{}'s configuration loader is closed. The documented contract separates
guidance from enforcement: CLAUDE.md files shape behavior but ``are not a
hard enforcement layer'', while settings rules are enforced by the client
regardless of model behavior; managed-policy memory files are
non-excludable and load first~\citep{claudeCodeMemoryDocs}. MCP server
definitions follow a five-level scope precedence (local, project, user,
plugin, connector) in which duplicate entries are taken wholesale from the
highest-precedence source without field merging, and project-scoped
\path{.mcp.json} servers require approval in interactive sessions and are
gated on workspace trust~\citep{claudeCodeMcpDocs}.

\subsection{State, rollout, and interfaces}
\label{sec:state}

\codex{} persists each thread as append-only JSONL under
\path{sessions/YYYY/MM/DD/}, flushing per record and retrying failed writes
once, with a \path{session_meta} head record; six SQLite databases mirror
thread metadata in WAL mode with a five-second busy timeout and are treated
as repairable, with filesystem read-repair fallback; and cold rollouts are
zstd-compressed after seven days~\citep{codexRepo}. Resume reconstructs the
transcript by reverse-scanning to a compaction base and replaying forward,
honoring rollbacks, which makes resume exact across compaction boundaries.
Memories are generated offline from rollouts by a two-phase pipeline.

\occ{} stores sessions in SQLite tables written exclusively by event
projectors over a durable event log with per-aggregate contiguous sequence
numbers, and separates file history from conversation history: snapshots
are git tree objects in a shadow repository (never commits), captured
before the stream and at every step boundary, capped by a 2 MiB untracked
file limit, and reverted by checking out saved hashes in batches of 100
files~\citep{opencodeRepo}. Sharing syncs sessions to a public URL with a
1000 ms debounce. Non-git projects get no snapshot tracking.

\ccc{} session storage is not inspectable. The documentation attests
per-subagent transcript files under the project directory, retained across
restarts within a session and deleted after a cleanup period that defaults
to 30 days, and hook inputs carry a \path{transcript_path} documented as
written asynchronously and possibly lagging the in-memory
conversation~\citep{claudeCodeSubagentsDocs,claudeCodeHooksDocs}.

Interfaces follow the same pattern of convergence with different
substrates. \codex{} dispatches TUI, exec, MCP server, app server, and
exec server from one binary, and every frontend is a client of one JSON-RPC
protocol in which approvals are server-to-client requests, backpressure has
a dedicated error code (-32001), and headless exec refuses to run outside a
git repository unless \path{--skip-git-repo-check} or the documented
bypass is set (\path{exec/src/lib.rs:799-805})~\citep{codexRepo}. \occ{}
runs the TUI against a local HTTP server exposing the full surface
(prompt, summarize, fork, abort, share, revert), implements the Agent
Client Protocol with one authentication method, and wires 38 built-in LSP
servers whose severity-one diagnostics are fed back into edit tool
outputs with the prefix ``LSP errors detected in this file, please
fix:''~\citep{opencodeRepo}. \ccc{}'s documentation attests CLI, IDE
extension, and desktop surfaces and a command-line MCP integration with
four transports; the reference deployment for untrusted work is a
firewalled devcontainer~\citep{claudeCodePermissionsDocs,%
claudeCodeMcpDocs,claudeCodeSurface}.

\subsection{Convergence and divergence}
\label{sec:convergence}

The evidence points to one shared architecture with three engineering
cultures.

\paragraph{Convergence.} (1) All three implement the abstract
prompt-sample-execute-continue loop and all three encode an explicit
anti-loop defense (rejection taxonomy; identical-call tripwire; stop-hook
cap). (2) The patch format is converging across vendors: \occ{} ports the
\codex{} \path{apply_patch} format verbatim in type names and selects it
for OpenAI model IDs, evidence of cross-harness borrowing at the protocol
level. (3) Extension vocabularies are converging on the \ccc{} design:
\codex{} exports \path{CLAUDE_PLUGIN_ROOT} compatibility variables and
discovers \path{.claude-plugin} manifests, and \occ{} reads \path{.claude}
skill layouts and CLAUDE.md files as fallbacks. (4) MCP is the shared
external-tool boundary in all three, with \codex{} exposing itself as an
MCP server (\path{codex} and \path{codex-reply} tools), and both open
systems naming MCP tools by server-prefixed identifiers. (5) All three
gate dangerous defaults behind enterprise layers (managed settings,
\path{requirements.toml}, MDM plists).

\paragraph{Genuine divergence.} (1) Control flow: task state machine
against persistence-driven loop against an undisclosed closed loop.
(2) Tool registration: conditional per-turn registration against a fixed
array against a partly documented inventory. (3) Safety substrate: OS
sandbox plus rule engine against in-process rules only against a Bash-scoped
OS sandbox plus permission layer. (4) Compaction triggers: fractional
window against reserved buffer against undisclosed. (5) State: JSONL-first
with a repairable SQL mirror against SQL-event-first with shadow-git file
history against undisclosed. (6) Interface philosophy: one protocol for
every frontend against a local HTTP server consumed by every frontend
against documented-only surfaces. Naming differences (hooks, skills,
agents/subagents, plugins) overlay these structural differences rather than
masking them.

\section{Related Work}
\label{sec:related}

Organized by comparison dimension.

\paragraph{Loop control and tool protocols.} ReAct established interleaved
reasoning and acting as the canonical agent loop and measured its
hallucination and repetition failure modes~\citep{yao2023react}; Toolformer
learned the tool-call emission protocol and constrained decoding to avoid
tool-call loops~\citep{schick2023toolformer}. The studied harnesses move
both concerns from prompting into machinery: typed rejection taxonomies,
identical-call tripwires, and stop-hook caps are engineering versions of
the same defenses.

\paragraph{Agent-computer interfaces and editing.} SWE-agent showed
interface design moves coding-agent performance: a guarded edit command and
lint gating each contributed measurable gains, while an iterative search
interface hurt~\citep{yang2024sweagent}. Our tool-surface findings extend
this from benchmark scaffolds to shipping products: all three systems
converge on structured edit protocols (grammar-parsed patches, staged
fuzzy replacers, documented edit tools), consistent with SWE-agent's
evidence that free-form shell editing is the failure-prone baseline.

\paragraph{Safety for tool-using agents.} ReAct's ethics statement limited
its action space to benign websites precisely because acting on real
environments is dangerous~\citep{yao2023react}; SWE-agent used Docker
containers with no permission model~\citep{yang2024sweagent}. The three
harnesses represent the next stage: per-command policy engines with OS
sandboxing, in-process permission rulesets, and layered permission-plus-
sandbox models with documented protected paths. Within our scope, none of
the three pinned sources cross-references the academic literature in its
policy design; the connection above is ours, grounded in the shared
failure patterns each mechanism addresses.

\paragraph{Third-party teardowns.} Three teardowns of \ccc{} were reviewed
and held to contextual use only: an intercepted-log teardown reporting
system-prompt and tool-token sizes and small-model routing
(minusXTeardown~\citep{minusXTeardown}), a binary-analysis catalog of
feature flags, telemetry events, and tools~\citep{tenguDecoded}, and a
trace-based breakdown of extension layers whose memory-delivery shape is
cross-verified by official documentation~\citep{agiflowInternals}. None
carries a comparison cell; their assertions are unverified against any
primary source.

\paragraph{Distinction from nearest work.} The nearest verifiable work is
the SWE-agent interface study, which ablates interface choices for one
external agent on one benchmark~\citep{yang2024sweagent}. This study
instead traces three closed-loop products at the implementation level,
covers eight machinery dimensions, and treats a closed-source system with a
documentation-tier protocol. Its bounds are correspondingly different:
no behavior was executed, and every claim is version-pinned.

\section{Limitations}
\label{sec:limitations}

\paragraph{Version bounds.} Every code claim is bounded to the three pinned
commits; every documentation claim to snapshots fetched 2026-08-20 from
sites that are not version-pinned and can drift. Nothing generalizes to
newer releases, and the \occ{} pin is the \texttt{dev} branch of the
rewrite, not a released version.

\paragraph{No behavior was measured.} The study is static by design. No
agent was run, no API was called, and no code from the pinned repositories
was executed. Statements about defaults and triggers describe code paths
and documentation contracts, not observed runtime behavior. The validation
runs re-check anchors but are smoke-test plumbing, not evidence about the
systems.

\paragraph{Closed core.} The \ccc{} turn loop, tool registry, compaction
thresholds, session storage, permission-enforcement internals, and sandbox
mechanisms are not established by this study; all are documented contracts
or surface attestations. The pinned documentation set also lacks the
compaction-survival page referenced by the memory docs, so the survival
rules beyond root CLAUDE.md are incomplete. Documentation describes
intended behavior, not necessarily shipped behavior on every platform.

\paragraph{Teardown tier.} Third-party teardowns are unverified claims
about captures and binaries this study cannot inspect; they appear only as
hedged context (Section~\ref{sec:tools}). One blog assertion (the
memory-delivery shape) is corroborated by official documentation, and only
that overlap is treated as reinforced.

\paragraph{Open cells.} Where no pinned source states a value, the
dimensions above say so explicitly: \ccc{} compaction thresholds, loop
implementation, full tool registry, and provider plumbing; \occ{}
permission merge precedence and runtime enforcement of several schema
defaults; \codex{} remote model-catalog and enterprise cloud-bundle
contents, which are server-side state. These are evidence gaps, not
findings of absence, except where stated (for example, the absence of an OS
sandbox in the pinned \occ{} tree is a searched-scope result).

\paragraph{Negative scope for Codex.} The extracted \codex{} registry has
no dedicated read/glob/grep tools within the searched components; absence
is bounded to the extraction scope and the pinned commit.

\section{Conclusion}

Three production harnesses implement the same abstract coding-agent loop
and diverge at every level below it. \codex{} builds a typed state machine
over abortable tasks, wraps shell access in an OS sandbox plus a rule
engine, and persists threads as inspectable JSONL with an exact replay
path. \occ{} builds a persistence-driven loop over a fixed tool array,
keeps all policy in-process, and invests its safety budget in recoverable
state, per-step snapshots with revert, rather than prevention. \ccc{}
presents a documentation-defined contract over a closed core: a wide hook
lifecycle, a two-layer permission model with a Bash-only OS sandbox, and an
extension surface whose stop-hook escape hatches already implement
autonomous loops outside the core. For a reader choosing or building a
harness, the comparison indicates that the abstraction is settled (loop,
tools, compaction, permissions, extensions) while the engineering culture
is not:
trust boundary placement, state durability, and interface unification are
the axes on which these systems genuinely differ. All conclusions are
bounded to the pinned commits and documentation snapshots named in
Section~\ref{sec:methods}.

\bibliographystyle{plainnat}
\bibliography{refs}

\end{document}
